
\medskip
\phantomsection\label{cor:sga3-VIB-4.3}%
\noindent\textbf{Corollary 4.3.---}
\emph{Let \(G\) be a flat \(S\)-group locally of finite
presentation. Then the function}
\[
  s\longmapsto\dim G_s
\]
\emph{is locally constant on \(S\).}

This follows immediately from \SGARefLink{prop:sga3-VIB-4.2}{4.2}, since every flat morphism
locally of finite presentation is universally open
(EGA~IV\(_2\), 2.4.6).

\medskip
\phantomsection\label{cor:sga3-VIB-4.4}%
\noindent\textbf{Corollary 4.4.---}
\emph{Let \(G\) be an \(S\)-group locally of finite presentation
over \(S\) at the points of the identity section. Consider the
following conditions:}
\begin{enumerate}
\item[\textup{\SGAInlineRefTarget{cor:sga3-VIB-4.4:item:i}(i)}] \emph{\(G\) is smooth over \(S\) at the points
  of the identity section \textup{(cf.\ \hyperref[thm:sga3-VIB-3.10]{3.10})}.}
\item[\textup{\SGAInlineRefTarget{cor:sga3-VIB-4.4:item:ii}(ii)}]\phantomsection\label{cond:sga3-VIB-4.4-ii} \emph{For every \(s\in S\), \(G_s\) is smooth
  over \(\kappa(s)\), and the function}
  \[
    s\longmapsto\dim G_s
  \]
  \emph{is locally constant on \(S\).}
\item[\textup{(iii)}] \emph{For every \(s\in S\), \(G_s\) is smooth
  over \(\kappa(s)\), and there exists a neighborhood \(V\) of the
  identity section such that the restriction to \(V\) of the
  structural morphism \(\pi:G\to S\) is universally open.}
\item[\textup{\SGAInlineRefTarget{cor:sga3-VIB-4.4:item:iv}(iv)}] \emph{For every \(s\in S\), \(G_s^0\) is
  smooth over \(\kappa(s)\), and \(G^0\) is representable by an
  open group subscheme of \(G\), universally open over \(S\).}
\end{enumerate}
\emph{Then the following implications hold:}
\[
  \textup{(i)}\Longrightarrow\textup{(ii)}
  \Longleftrightarrow\textup{(iii)}
  \Longleftrightarrow\textup{(iv)}.
\]
\emph{If, moreover, \(S\) is reduced, then conditions
\textup{\SGARefLink{cor:sga3-VIB-4.4:item:i}{(i)}}--\textup{\SGARefLink{cor:sga3-VIB-4.4:item:iv}{(iv)}} are equivalent and imply that \(G^0\)
is smooth over \(S\).}\footnote{Translator's note: the Gille--Polo re-edition (2011),
Expos\'e~VI\textsubscript{B}, shows that \(S\) must be reduced here. In our formulation: let \(k\) be a
field, \(A=k[\delta]\) with \(\delta^2=0\), \(S=\operatorname{Spec}A\), and let \(G\) be the closed subgroup
\(\operatorname{Spec}A[X]/(\delta X)\) of \(\mathbb G_{a,S}\), whose points with values in an \(A\)-algebra
\(R\) are the \(r\in R\) with \(\delta r=0\). Its only fibre is \(\mathbb G_{a,k}\), smooth and connected;
\(G=G^0\) is universally open over \(S\), having the same underlying space as \(\mathbb G_{a,k}\) after
any base change; so \textup{(ii)}--\textup{(iv)} hold. But \(G\) is not flat over \(A\), since the
coefficient module \(A/\delta A\) of \(X\) is not flat over \(A\); so \textup{(i)} fails.}

Let us show that \textup{\SGARefLink{cor:sga3-VIB-4.4:item:i}{(i)}} implies \textup{\SGARefLink{cor:sga3-VIB-4.4:item:ii}{(ii)}}. By
EGA~IV\(_4\), 17.10.2, the function
\[
  x\longmapsto
  \dim_x\pi^{-1}\bigl(\pi(x)\bigr)
    =\dim\pi^{-1}\bigl(\pi(x)\bigr)
  \qquad(\hyperref[lem:sga3-VIB-1.5]{1.5})
\]
is continuous in a neighborhood of the identity section; hence the
function \(s\mapsto\dim G_s\) is continuous on \(S\), and therefore locally
constant on \(S\). On the other hand, for every \(s\in S\), \(G_s\) is
smooth over \(\kappa(s)\) (\SGARefLink{cor:sga3-VIB-1.3.1}{1.3.1}).

Let us show that \textup{\SGARefLink{cor:sga3-VIB-4.4:item:ii}{(ii)}} implies \textup{\SGARefLink{cor:sga3-VIB-4.4:item:iv}{(iv)}}. The question
is local on \(S\), since it suffices to show that \(\underline{G^0}\) is open
in \(G\) (\hyperref[case:sga3-VIB-3.4]{3.4}), the properties of \(G^0\) stated above then following
from \hyperref[cor:sga3-VIB-2.4]{2.4} and \hyperref[prop:sga3-VIB-4.2]{4.2}. Given \(s\in S\), let us construct, as in the proof
of \hyperref[thm:sga3-VIB-3.10]{3.10}, \(W\), \(T\), \(\pi''\), \(\epsilon''\) and \(W^0\). It then follows from
EGA~IV\(_3\), 15.6.7, that \(W^0\) is open in \(W\). It follows from
\hyperref[prop:sga3-VIB-4.2]{4.2} that \(\pi\) is universally open at every point of \(W^0\), so that
(Expos\'e~VI\(_A\), \SGARefLink{sga3:VIA:0.1}{0.1}) \(\mu\) is universally open at every point of
\(W^0\times_S W^0\), which shows that \(W^0\cdot W^0\) is open; and one concludes
as in the proof of \hyperref[thm:sga3-VIB-3.10]{3.10}.

It is clear that
\[
  \textup{(iv)}\Longrightarrow\textup{(iii)},
\]
and \(\textup{(iii)}\Longrightarrow\textup{(ii)}\) follows by
applying \SGARefLink{prop:sga3-VIB-4.2}{4.2} to \(V\).

Let us finally show that, if \(S\) is reduced, \textup{(iv)} implies \textup{(i)}: we may
evidently suppose that \(G=G^0\), so that \(G\) is of finite presentation over
\(S\) by \SGARefLink{cor:sga3-VIB-5.5}{5.5} below,\footnote{Translator's note: the 1970 text refers to
5.4; the statement used, for a group with connected fibers, universally open over \(S\), is 5.5, as in the
Gille--Polo re-edition (2011), Expos\'e~VI\textsubscript{B}. The re-edition also details the use of
EGA~IV\(_3\), 15.6.7: the fibers of \(G\) are smooth and connected, hence geometrically integral, and their
dimension is locally constant, so that \(G\) is flat over \(S_{\mathrm{red}}=S\).} and the conclusion then follows from
EGA~IV\(_3\), 15.6.7 (which implies that \(G\) is flat over \(S\)), and from
\hyperref[thm:sga3-VIB-3.10]{3.10}.
